\section{Theorem 41: the cut-graded universal generator}

\subsection{Unique cut grading}

Define
\[
 \Ge=\frac12(G+JGJ),
 \qquad
 \Go=\frac12(G-JGJ).
\]

\begin{theorem}[Theorem 41A: unique cut grading]\label{thm:t41grading}
The generator decomposes uniquely as
\[
 G=\Ge+\Go,
\]
with
\[
 J\Ge J=\Ge,
 \qquad
 J\Go J=-\Go.
\]
Moreover,
\[
 P_+\Ge P_-=P_-\Ge P_+=0,
\]
while
\[
 P_+\Go P_+=P_-\Go P_-=0.
\]
Thus $\Ge$ is sheet-preserving and $\Go$ is trans-cut.
\end{theorem}

\begin{proof}
The sum identity follows by addition. Conjugating by $J$ and using $J^2=I$ gives the two parity relations. If $G=A+B$ with $JAJ=A$ and $JBJ=-B$, then
\[
 A=\frac12(G+JGJ),\qquad B=\frac12(G-JGJ),
\]
so the decomposition is unique. Multiplication by $P_\pm$ yields the block-support statements.
\end{proof}

\subsection{Bilateral reconstruction and the operator cut square}

Assume now that $G$ is cut odd:
\[
 JGJ=-G.
\]

\begin{theorem}[Theorem 41B: bilateral exponential reconstruction]\label{thm:t41bilateral}
For every $t$ for which the exponential is defined,
\[
 JU_tJ=U_{-t}.
\]
Put
\[
 E_t=\frac12(U_t+U_{-t}),
 \qquad
 O_t=\frac12(U_t-U_{-t}).
\]
Then
\[
 U_t=E_t+O_t,
 \qquad
 U_{-t}=E_t-O_t,
\]
\[
 JE_tJ=E_t,
 \qquad
 JO_tJ=-O_t,
\]
and
\[
 E_t^2-O_t^2=I.
\]
Equivalently, with
\[
 \widehat J_t=U_t+U_{-t},
 \qquad
 \widehat C_t=U_t-U_{-t},
\]
one has the operator cut-square identity
\[
 \boxed{\widehat J_t^2-\widehat C_t^2=4I.}
\]
The two channels commute and
\[
 \widehat J_t\widehat C_t
 =\widehat C_t\widehat J_t
 =U_{2t}-U_{-2t}.
\]
\end{theorem}

\begin{proof}
Functional calculus and $JGJ=-G$ give
\[
 Je^{tG}J=e^{tJGJ}=e^{-tG}.
\]
The parity and reconstruction formulas follow by addition and subtraction. Because $U_t$ and $U_{-t}$ are functions of the same operator, they commute and $U_tU_{-t}=I$. Expanding gives
\[
 E_t^2-O_t^2
 =\frac14\bigl[(U_t+U_{-t})^2-(U_t-U_{-t})^2\bigr]=I.
\]
Multiplying by four gives the boxed identity; direct expansion gives the product formula.
\end{proof}

\begin{remark}
The operator cut square is a bilateral flow identity. It is not the positive-metric Lagrange identity proved in \cref{sec:metric-cut-square}. Both are called cut-square identities because each separates a recognised channel from its complementary structure, but their carriers, signs, and hypotheses differ.
\end{remark}

\subsection{Cut-loop memory and seam curvature}

For a general bounded generator define
\[
 \mathscr H_J(t)=JU_tJU_t=e^{tJGJ}e^{tG}.
\]

\begin{theorem}[Theorem 41C: logarithmic cut-loop expansion]\label{thm:t41loop}
As $t\to0$,
\[
 \mathscr H_J(t)
 =I+2t\Ge+t^2\bigl(2\Ge^2+[\Ge,\Go]\bigr)+O(t^3),
\]
and
\[
 \boxed{
 \log\mathscr H_J(t)
 =2t\Ge+t^2[\Ge,\Go]+O(t^3).}
\]
Consequently,
\[
 \Ge=\frac12\left.\frac{d}{dt}\right|_{t=0}\mathscr H_J(t),
 \qquad
 \mathcal R_{\mathrm{seam}}=[\Ge,\Go].
\]
The cut loop is locally the identity exactly when $\Ge=0$.
\end{theorem}

\begin{proof}
Substitute $JGJ=\Ge-\Go$ and $G=\Ge+\Go$ into the product of exponential series. The linear term is $2\Ge$. The quadratic term is $2\Ge^2+[\Ge,\Go]$. For $I+tA+t^2B+O(t^3)$,
\[
 \log(I+tA+t^2B+O(t^3))=tA+t^2\left(B-\frac12A^2\right)+O(t^3),
\]
which yields the logarithmic formula. If the loop is locally $I$, its first derivative vanishes, hence $\Ge=0$. Conversely, if $\Ge=0$, then $JGJ=-G$ and $JU_tJ=U_{-t}$, so the loop equals $I$.
\end{proof}

\subsection{Closure sectors and the clock-free Eye}

\begin{proposition}[Periodic and antiperiodic sectors]
If $G\psi=i\omega\psi$ with $\omega\in\R$, then for integer $N\ge1$,
\[
 U_N\psi=\psi \iff N\omega\in2\pi\mathbb Z,
\]
while
\[
 U_N\psi=-\psi \iff N\omega\in(2\mathbb Z+1)\pi.
\]
For $\sigma\in\{\pm1\}$,
\[
 S_{N,\sigma}=(U_N-\sigma I)^*(U_N-\sigma I)\ge0,
\]
and
\[
 \Ker S_{N,\sigma}=\Ker(U_N-\sigma I).
\]
\end{proposition}

\begin{proof}
Apply the exponential to the eigenvector. Positivity and the kernel identity follow from the definition of $S_{N,\sigma}$.
\end{proof}

Assume now $U_t=e^{-itH}$ with $H=H^*$. The one-step fixed space can contain nonzero integer frequencies:
\[
 \Fix(U_1)=E_H(2\pi\mathbb Z)\Hh.
\]

\begin{theorem}[Theorem 41D: clock-free Eye]\label{thm:eye}
The all-time fixed space is
\[
 \bigcap_{t\in\R}\Fix(U_t)=\Ker H,
\]
and the mean-ergodic projector is
\[
 \boxed{
 \Eye=\operatorname*{s-lim}_{T\to\infty}
 \frac1T\int_0^TU_t\,dt=E_H(\{0\}).}
\]
Thus
\[
 \Fix(U_1)/\Ker H
\]
is the stroboscopic recognition-blind sector.
\end{theorem}

\begin{proof}
The spectral theorem gives
\[
 \frac1T\int_0^TU_t x\,dt
 =\int_{\R}\left(\frac1T\int_0^Te^{-it\lambda}\,dt\right)dE_H(\lambda)x.
\]
The scalar factor tends to one at $\lambda=0$ and zero elsewhere, and is bounded in modulus by one. Dominated convergence with respect to the spectral measure proves the strong limit.
\end{proof}

\subsection{Faithful component observer}

Let $P_1,\ldots,P_m$ be pairwise orthogonal projections with $\sum_aP_a=I$. Define $G_a=P_aG$ and
\[
 \Aobs_Gx=(G_1x,\ldots,G_mx).
\]

\begin{theorem}[Theorem 41E: component Gram preservation]
One has
\[
 \Aobs_G^*\Aobs_G=\sum_aG_a^*G_a=G^*G,
\]
and therefore
\[
 \Ker\Aobs_G=\Ker G.
\]
\end{theorem}

\begin{proof}
\[
 \sum_aG_a^*G_a=G^*\left(\sum_aP_a\right)G=G^*G.
\]
The kernel equality follows because $\|\Aobs_Gx\|^2=\|Gx\|^2$.
\end{proof}

This result is load-bearing for two-axis nuclear observation: a direct sum preserves both directional energies, whereas a scalar aggregate can vanish by cancellation.
